% !TeX root = ../../Book5.tex
% 纲要标题：### B.4 参数特殊化与理论的适用范围
% 纲要位置：工作记录/计划纲要.md，第 4601 行。
% #### B.4.1 单位根处的表示变化
% #### B.4.2 有理变异与点上的退化
\section{参数特殊化与理论的适用范围}
\label{b5:sec:B04}

量子参数和簇变量都能代入具体值，但代入可能使原来可逆的系数消失。要比较一般参数与特殊参数，须先确定在哪一个整系数或 Laurent 系数环上定义对象，再说明取商后哪些结论仍然成立。

\subsection{单位根处的表示变化}
\label{b5:subsec:B04:01}

在 \(U_q(\mathfrak{sl}_2)\) 的模型中，不可约性证明使用
\([i]_q\ne0\) 以及 \(K\) 的权互异。若 \(q\) 为奇数阶 \(\ell\ge3\) 的本原单位根，则 \([\ell]_q=0\)。例如取 \(m=\ell\)，向量
\(v_1,\ldots,v_\ell\)
生成一个非零真子模：\(K\) 保持它，\(F\) 只向后移，且
\(Ev_1=[1]_q[\ell]_qv_0=0\)，其余 \(Ev_i\) 仍在其中。所以一般参数下的不可约性在这个值上失效。
\ProseParagraphBreak
把 \(q\) 取为一时，不能直接代入
\((K-K^{-1})/(q-q^{-1})\)，因为分母变为零。须先选择适当的整形式或引入额外生成元，再讨论特殊化。不同整形式在单位根处可能给出不同代数，见\cite[Section 5.6, Remark 5.6.5]{EtingofEtAl2015TensorCategories}；形式极限本身不足以确定特殊化后的代数。

\subsection{有理变异与点上的退化}
\label{b5:subsec:B04:02}

簇变异是一项可逆有理坐标变换，不是每个仿射点上都定义的正则同构。前面的 \(A_2\) 例子中，若取 \(x=y=-1\)，五个 Laurent 表达均有意义，得到
\[
(z_1,z_2,z_3,z_4,z_5)=(-1,-1,0,-1,0).
\]
交换恒等式仍成立，但若从含 \(z_3=0\) 的簇出发，用除以 \(z_3\) 的公式恢复邻近变量，就不能直接在这个点上操作。一个在有理函数域中合法的逆变换，可能在特定点失去意义。
\ProseParagraphBreak
计算时应区分这两种情形：变量作为代数独立有理函数时，变异可逆且满足已证恒等式；变量代入一个点时，需要逐项检查所用分母。类似的区分也适用于量子参数的特殊化和范畴化中的基变换。进一步学习时，还需分别讨论规范基、正性、量子簇结构和高阶范畴化的定义与定理。
